\section{真机实验}
\label{app:real}


\subsection{平台与任务}
\label{app:real_platform}
机器人为\cfgRobotZh{}。策略以 20\,Hz 下发工具中心点的绝对笛卡尔位姿与夹爪目标。采集示范时，以 100\,Hz 记录机器人控制数据。平台配有两台 RealSense D435，分别安装在固定支架和腕部，均以 $640\times480$、60\,Hz 采集 RGB。我们用 ChArUco 标定板求固定相机到机器人基座的变换。正式评测计划为 \ours 提供固定相机的 RGB-D 帧及该视角估计的目标位姿，为基线提供两台相机的 RGB。各任务装置轮流安装在这一平台上（图~\mainref{fig:real_trials}；表~\ref{tab:real_tasks}）。


\begin{table}[!htbp]
\centering\footnotesize
\caption{\textbf{真机任务。}稳定夹住目标并将其抬离支撑面即为抓取成功。所有任务的试次时限均为 30\,s。}
\label{tab:real_tasks}
\setlength{\tabcolsep}{4pt}
\begin{tabularx}{\linewidth}{@{}l>{\raggedright\arraybackslash}Xl@{}}
\toprule
任务 & 目标运动 & 成功判据 \\
\midrule
T1 & 方块随传送带移动 & 抓取 \\
T2 & 无盖盒随传送带移动 & 方块留在盒内 \\
T3 & 长方块随转盘运动 & 沿短边抓取 \\
T4 & 玩具车在桌面上移动 & 抓取 \\
T5 & 网球滚动，初速度大小和起始方向偏移随机化 & 夹爪使球停下或转向 \\
T6 & 斜坡小车载着方块移动 & 抓取 \\
\bottomrule
\end{tabularx}
\end{table}

\subsection{位姿跟踪与标签}
\label{app:tracker}
\paragraph{离线标签。}
我们对固定相机的录像离线标注。Grounding DINO \cite{groundingdino} 与 SAM~2.1 \cite{sam2} 生成目标掩码，FoundationPose \cite{foundationpose} 利用 FoundationStereo \cite{foundationstereo} 深度和米制物体网格做正反向跟踪。我们核对两向跟踪的一致性，并用工具中心点的运动学补齐刚性夹持时受遮挡的帧。位姿标签保留 595 条示范，其中 97.1\% 的帧有有效目标位姿。网球只标注平移。稀疏场景流将 CoTracker3 \cite{cotracker3} 的点轨迹与双目深度结合，映射到机器人基座坐标系。我们对有效观测插值，取得 50--400\,ms 各预测时域的监督。

\paragraph{在线位姿估计。}
我们从固定相机的对齐 RGB-D 帧估计目标位姿，并用相机外参将其转换到机器人基座坐标系。策略以 20\,Hz 采样位姿估计，构建目标状态历史。


\subsection{训练与评测协议}
\label{app:real_trials}
共享数据集含 600 条遥操作示范，T1 至 T6 分别有 100、100、101、100、99 和 100 条。ACT 为每个任务单独训练 50{,}000 步。DynamicWAM 在六个任务上联合微调 20{,}000 步。其部署配置采用 16 步异步动作块，旧队列还剩 4 步（200\,ms）时查询新动作块，并融合重叠动作。\ours 的微调见附录~\ref{app:posttraining}。

试次按任务和方法分组，每组 20 次。操作员在每轮开始前重置场景，并在结束后记录结果。超时、人工干预和安全中止均计为失败。表~\ref{tab:supp_timing} 给出真机计时。

\subsection{结果与失败分析}
\label{app:real_results}
% Planned results: six-task real-robot study (App. D.3); per-task counts, baseline set, and numeric results to be filled (see docs/writing/v4_t5_real_report.md).
\begin{table*}[t]
\centering\footnotesize
\caption{\tablelang{\textbf{Real-robot successes} per task (successes/trials, \fillin{n} trials per task; tasks in Table~\ref{tab:real_tasks}). All: pooled success rate (\%) with its 95\% binomial interval. $\Delta$: paired difference from \ours (points) with a bootstrap interval over trial keys. $e_t$: median contact-time error (ms) on trials that make contact.}{\textbf{真机成功次数}，按任务分列（成功次数/试验数，每个任务 \fillin{n} 次；任务定义见表~\ref{tab:real_tasks}）。All：合并成功率（\%）及 95\% 二项区间。$\Delta$：相对 \ours 的配对差值（百分点）及按试验键自助法得到的区间。$e_t$：发生接触的试验中接触时刻误差的中位数（ms）。}}
\label{tab:supp_real_ci}
\setlength{\tabcolsep}{3pt}
\begin{tabular}{@{}lccccccccc@{}}
\toprule
\tablelang{Method}{方法} & T1 & T2 & T3 & T4 & T5 & T6 & All & $\Delta$ & $e_t$ \\
\midrule
\fillin{$\pi_{0.5}$} & \resultplaceholder[s/n] & \resultplaceholder[s/n] & \resultplaceholder[s/n] & \resultplaceholder[s/n] & \resultplaceholder[s/n] & \resultplaceholder[s/n] & \resultplaceholder[00.0 {[00, 00]}] & \resultplaceholder[{$-$00.0 [00, 00]}] & \resultplaceholder \\
\fillin{DynamicWAM} & \resultplaceholder[s/n] & \resultplaceholder[s/n] & \resultplaceholder[s/n] & \resultplaceholder[s/n] & \resultplaceholder[s/n] & \resultplaceholder[s/n] & \resultplaceholder[00.0 {[00, 00]}] & \resultplaceholder[{$-$00.0 [00, 00]}] & \resultplaceholder \\
\tablelang{\ours, no DROID base}{\ours，无 DROID 基座} & \resultplaceholder[s/n] & \resultplaceholder[s/n] & \resultplaceholder[s/n] & \resultplaceholder[s/n] & \resultplaceholder[s/n] & \resultplaceholder[s/n] & \resultplaceholder[00.0 {[00, 00]}] & \resultplaceholder[{$-$00.0 [00, 00]}] & \resultplaceholder \\
\ours & \resultplaceholder[s/n] & \resultplaceholder[s/n] & \resultplaceholder[s/n] & \resultplaceholder[s/n] & \resultplaceholder[s/n] & \resultplaceholder[s/n] & \resultplaceholder[00.0 {[00, 00]}] & -- & \resultplaceholder \\
\bottomrule
\end{tabular}
\end{table*}

\begin{figure*}[!htbp]
\centering
\figplaceholder[2.1in]{\textwidth}{真机失败}{%
\textbf{(a) 失败类别。}每个方法用一根堆叠条显示各类失败次数。
\textbf{(b) 代表性试次。}展示未经修饰的固定视角画面，标出决定结果的事件。}
\caption{\textbf{真机失败。}(a) 各方法按类别统计的失败次数。(b) 代表性试次的画面。\fillin{所示试次的任务、方法和结果。}}
\label{fig:real_failures}
\end{figure*}
表~\ref{tab:supp_real_ci} 给出合并成功率及其区间。图~\ref{fig:real_failures} 按可观察事件归类失败，包括抓取过晚或未对准、滑落或掉落、碰撞或安全中止、超时，以及目标跟踪丢失。\fillin{各方法的主要失败类型和判断时序误差的依据。}
